\label{sec:XML:AppLocalizability}

How much of what a production-style detector relies on can be placed in the input at all? The null element makes this measurable: its share of a pulled-back explanation is exactly the attribution that no part of the input can carry (Definition~\ref{def:Pullback}). We report it for the explanation set (2{,}000 files, one seed).

In a typical file, about 70\% of the 2{,}568 features are unattributable, almost all because something is absent (e.g., an import hash bucket that nothing falls into); 102 features are never attributable, 88 of them because the reference extractor does not report which bytes caused a parser warning. About 22\% of the GBDT's attribution (TreeSHAP) rests on unattributable features, 17--18\% of it on absences; for the MLP (IG), 37\% rests on absences alone. Explainers whose baseline is zero (gradient $\times$ input) cannot credit absences at all, so the measurement must be reported per explainer.
